\documentclass[11pt, a4paper]{article}

% --- UNIVERSAL PREAMBLE BLOCK ---
\usepackage[a4paper, top=2.5cm, bottom=2.5cm, left=2cm, right=2cm]{geometry}
\usepackage{fontspec}
\usepackage[english, bidi=basic, provide=*]{babel}

\babelprovide[import, onchar=ids fonts]{english}

% Set default/Latin font to Sans Serif in the main (rm) slot
\babelfont{rm}{Noto Sans}

\usepackage{amsmath}
\usepackage{amssymb}
\usepackage{amsthm}
\usepackage{booktabs}
\usepackage{enumitem}
\setlist[itemize]{label=-}

\usepackage[colorlinks=true, linkcolor=blue, urlcolor=blue, citecolor=blue]{hyperref}

% Theorem environments
\newtheorem{definition}{Definition}
\newtheorem{theorem}{Theorem}
\newtheorem{lemma}{Lemma}
\newtheorem{corollary}{Corollary}
\newtheorem{proposition}{Proposition}
\newtheorem{remark}{Remark}

\title{The COD Kernel Representation Framework: \\ Operator-Algebraic Reconstruction of Quantum Correlation Geometry}
\author{Omega Protocol Formalization}
\date{April 2026}

\begin{document}

\maketitle

\begin{abstract}
We present an operator-algebraic framework that reconstructs aspects of quantum correlation geometry from Completely Positive Trace-Preserving (CPTP) informational flow. By explicitly separating functional analytic derivations from physical interpretations, we show how non-commutative probability structures can be induced from families of Choi operators. We construct a positive-semidefinite Gram kernel via the Hilbert--Schmidt inner product of Choi states and use it to define a positive functional on the operator system generated by node observables. Using the Gelfand--Naimark--Segal (GNS) construction, we obtain a cyclic representation in which the kernel's imaginary part controls commutator expectation values. For bipartite settings, we recover Tsirelson-type bounds with standard operator-norm assumptions. Finally, we state an explicit asymptotic-commutativity hypothesis under which the represented algebra abelianizes in the large-time limit.
\end{abstract}

\section{Introduction}

The objective of this framework is to establish a mathematically rigorous pipeline from informational operations to quantum geometric structures. To avoid circularity, this text strictly delineates rigorous functional analysis (Theorem and Definition layers) from physical correspondences (Interpretation layers).

\section{The Gram Kernel Space}

Let $\mathcal{N}$ be a discrete index set of interacting nodes. For each $i \in \mathcal{N}$, let $\mathcal{E}_i$ be a Completely Positive Trace-Preserving (CPTP) map with Choi--Jamio\l{}kowski operator $J_i$.

\begin{definition}[Informational Gram Kernel]
Assume $J_i$ are Hilbert--Schmidt operators on a common Hilbert space. Define
\begin{equation}
    K(i,j) = \langle J_i, J_j \rangle_{HS} = \mathrm{Tr}(J_i^\dagger J_j),
\end{equation}
for $i,j \in \mathcal{N}$.
\end{definition}

\begin{proposition}[Canonical Hermitian Decomposition]
As a Gram kernel over a complex Hilbert space, $K$ is Hermitian ($K(i,j)=\overline{K(j,i)}$) and positive semidefinite. Hence
\begin{equation}
    K(i,j) = g_{ij} + i\,\omega_{ij},
\end{equation}
where $g_{ij}=\operatorname{Re}K(i,j)$ is symmetric and $\omega_{ij}=\operatorname{Im}K(i,j)$ is skew-symmetric.
\end{proposition}

\begin{remark}[Mathematical Constraint]
The positivity condition $K \succeq 0$ constrains the admissible pair $(g,\omega)$. In finite dimensions, this is equivalent to positivity of the complex matrix $g+i\omega$.
\end{remark}

\begin{remark}[Physical Interpretation]
We identify $g_{ij}$ as a Chain Overlap Density (COD), representing metric overlap, and $\omega_{ij}$ as a Reverse Chain Overlap Density (RCOD), representing directional causal flux.
\end{remark}

\section{Emergence of Observables via GNS Representation}

\begin{definition}[Operator-System State Induced by the Kernel]
Let $\mathcal{V}=\operatorname{span}\{1,A_i,A_i^\dagger, A_i^\dagger A_j\,:\, i,j\in\mathcal{N}\}$ be the operator system generated up to degree two. Define
\begin{equation}
\omega_{\mathrm{state}}(A_i^\dagger A_j):=K(i,j), \qquad \omega_{\mathrm{state}}(1):=1,
\end{equation}
and extend linearly on $\mathcal{V}$. Positivity on degree-two elements follows from $K\succeq 0$.
\end{definition}

\begin{theorem}[GNS Emergence (after positive extension)]
If $\omega_{\mathrm{state}}$ admits a positive extension to a unital $C^*$-algebra $\mathcal{A}$ containing $\mathcal{V}$, then by the GNS construction there is a cyclic triple $(\pi,\mathcal{H}_{\mathrm{GNS}},\Omega)$ with
\begin{equation}
    \omega_{\mathrm{state}}(X)=\langle \Omega,\pi(X)\Omega\rangle, \quad X\in\mathcal{A},
\end{equation}
and in particular
\begin{equation}
    K(i,j)=\langle \Omega,\pi(A_i)^\dagger\pi(A_j)\Omega\rangle.
\end{equation}
\end{theorem}

\section{Symplectic Form and the Central Extension}

\begin{proposition}[Expectation-Value 2-Form]
For represented observables $\hat E_i:=\pi(A_i)$,
\begin{equation}
    \langle\Omega,[\hat E_i,\hat E_j]\Omega\rangle
    =\langle\Omega,\hat E_i\hat E_j\Omega\rangle-\langle\Omega,\hat E_j\hat E_i\Omega\rangle
    =2i\,\operatorname{Im}\langle\Omega,\hat E_i\hat E_j\Omega\rangle.
\end{equation}
If $\langle\Omega,\hat E_i\hat E_j\Omega\rangle=K(i,j)$ (e.g., for self-adjoint generators with the chosen state), then
\begin{equation}
\langle\Omega,[\hat E_i,\hat E_j]\Omega\rangle=2i\,\omega_{ij}.
\end{equation}
\end{proposition}

\begin{remark}[Operator Identity vs. Expectation Identity]
The scalar $\omega_{ij}$ is an expectation-level 2-form. It does not by itself imply an operator identity $[\hat E_i,\hat E_j]=2i\omega_{ij}I$. Such identities require additional algebraic assumptions (e.g., a CCR/Weyl representation and centrality conditions).
\end{remark}

\section{Bipartite Embedding and Correlation Bounds}

\begin{theorem}[Tsirelson-Type Bound in the Bipartite Representation]
Consider a bipartite representation with observables $A_0,A_1$ on subsystem $A$ and $B_0,B_1$ on subsystem $B$, each self-adjoint and contractive ($\|A_x\|\le 1$, $\|B_y\|\le 1$). For the CHSH operator
\begin{equation}
\mathcal{B}:=A_0\otimes(B_0+B_1)+A_1\otimes(B_0-B_1),
\end{equation}
one has $\|\mathcal{B}\|\le 2\sqrt{2}$. Therefore every state $\rho$ satisfies
\begin{equation}
|\operatorname{Tr}(\rho\,\mathcal{B})|\le 2\sqrt{2}.
\end{equation}
\end{theorem}

\begin{remark}[Physical Interpretation]
This yields a Tsirelson bound from operator-norm constraints in the embedded bipartite representation, without requiring a separate axiomatization beyond the representation framework.
\end{remark}

\section{Asymptotic Abelianization and Classicality}

\begin{proposition}[Asymptotic Commutativity Hypothesis]
Let $t\mapsto A_i^{(t)}$ be evolved observables in the represented algebra. Assume
\begin{equation}
\lim_{t\to\infty}\|[\pi(A_i^{(t)}),\pi(A_j^{(t)})]\|=0\quad\text{for all }i,j.
\end{equation}
Then commutators vanish in every vector state in the limit, so the limiting correlation structure is abelian at the level of observables/correlators. If, additionally, the von Neumann algebra generated by the limit observables exists, it is commutative.
\end{proposition}

\begin{remark}[Scope]
The commutator-norm decay above is an explicit assumption (or theorem premise), not a consequence of CPTP dynamics alone. Any application should verify it for the concrete dynamical model used.
\end{remark}

\section{Conclusion}

Separating Gram-kernel data from operator-representation data yields a clean chain:
\[
\text{CPTP maps}\to \text{Gram kernel}\to \text{positive functional (extension)}\to \text{GNS representation}\to \text{operator correlations}.
\]
Within this chain, the imaginary kernel component controls commutator expectations, Tsirelson-type bounds follow from bipartite contractive observables, and classical behavior appears under explicit asymptotic-commutativity hypotheses.

\begin{thebibliography}{9}
\bibitem{GNS} Gelfand, I. \& Naimark, M. (1943). On the imbedding of normed rings into the ring of operators in Hilbert space. \textit{Recueil Math\'ematique}, 12(54).
\bibitem{QuantumGrothendieck} Pisier, G. \& Shlyakhtenko, D. (2002). Grothendieck's theorem for operator spaces. \textit{Inventiones Mathematicae}, 150(1).
\bibitem{ExtremeEvents} Papastathopoulos, I. et al. (2016). Extreme Events of Markov Chains. \textit{arXiv:1510.08920}.
\end{thebibliography}

\end{document}
